\documentclass[
    11pt,
    a4paper,
    oneside
]{article}

% ---------- Language and typography ----------
\usepackage[T1]{fontenc}
\usepackage[utf8]{inputenc} % Omit when using LuaLaTeX/XeLaTeX
\usepackage[english]{babel}
\usepackage{lmodern}
\usepackage{microtype}

% ---------- Page layout ----------
\usepackage[
    top=25mm,
    bottom=25mm,
    left=30mm,
    right=25mm
]{geometry}
\usepackage{setspace}
\onehalfspacing
\usepackage{parskip}
\setlength{\parindent}{0pt}

% ---------- Mathematics ----------
\usepackage{amsmath,amssymb,mathtools}
\usepackage{bm}
\usepackage{siunitx}
\sisetup{
    separate-uncertainty=true,
    per-mode=symbol,
    range-phrase=--,
    range-units=single
}

% Useful mathematical operators
\DeclareMathOperator{\diag}{diag}
\DeclareMathOperator{\RePart}{Re}
\DeclareMathOperator{\ImPart}{Im}
\newcommand{\dd}{\mathop{}\!\mathrm{d}}
\newcommand{\e}{\mathrm{e}}
\newcommand{\iu}{\mathrm{i}}
\newcommand{\vect}[1]{\bm{#1}}
\usepackage{amsmath,amssymb,mathtools}


\usepackage{amsmath}
\usepackage{graphicx}
\usepackage[europeanresistors]{circuitikz}
\ctikzset{bipoles/length=1.3cm}

% Circuit notation
\newcommand{\seriesconn}{\mathbin{\oplus}}
\newcommand{\Lel}[1]{\mathrm{L}\!\left(#1\right)}
\newcommand{\Wel}[1]{\mathrm{W}\!\left(#1\right)}
\newcommand{\Gel}[2]{\mathrm{G}\!\left(#1,#2\right)}
\newcommand{\CPEel}[2]{\mathrm{CPE}\!\left(#1,#2\right)}


% Grammar notation
\newcommand{\nt}[1]{\mathit{#1}}
\newcommand{\term}[1]{\text{\texttt{#1}}}
\newcommand{\defines}{\Coloneqq}
% Fractional derivatives
\newcommand{\Caputo}[2]{%
    {}^{\mathrm{C}}\!D_{#1}^{#2}%
}

% ---------- Figures and tables ----------
\usepackage{graphicx}
\graphicspath{{figures/}}
\newcommand{\reportgraphic}[2][]{%
    \IfFileExists{#2}{%
        \includegraphics[#1]{#2}%
    }{%
        \IfFileExists{figures/#2}{%
            \includegraphics[#1]{figures/#2}%
        }{%
            \PackageWarning{report}{Image `#2' not found; using a placeholder}%
            \fbox{%
                \parbox[c][45mm][c]{0.9\linewidth}{%
                    \centering Missing image: \texttt{\detokenize{#2}}%
                }%
            }%
        }%
    }%
}
\usepackage{booktabs}
\usepackage{tabularx}
\usepackage{multirow}
\usepackage{array}
\usepackage{subcaption}
\usepackage{float}

% ---------- Algorithms and source code ----------
\usepackage{algorithm}
\usepackage{algpseudocode}
\usepackage{listings}
\usepackage{xcolor}

\definecolor{codegreen}{RGB}{0,120,70}
\definecolor{codegray}{RGB}{100,100,100}
\definecolor{codeblue}{RGB}{20,70,160}
\definecolor{codebackground}{RGB}{247,247,247}

\lstset{
    backgroundcolor=\color{codebackground},
    basicstyle=\ttfamily\small,
    keywordstyle=\color{codeblue}\bfseries,
    commentstyle=\color{codegreen},
    stringstyle=\color{red!70!black},
    numbers=left,
    numberstyle=\tiny\color{codegray},
    frame=single,
    breaklines=true,
    showstringspaces=false,
    tabsize=4,
    captionpos=b
}

% ---------- References ----------
\usepackage[
    backend=biber,
    style=numeric-comp,
    sorting=none,
    giveninits=true,
    doi=true,
    url=true
]{biblatex}
\addbibresource{references.bib}

% ---------- Cross-references and links ----------
\usepackage{csquotes}
\usepackage{hyperref}
\hypersetup{
    colorlinks=true,
    linkcolor=blue!50!black,
    citecolor=green!40!black,
    urlcolor=blue!60!black,
    pdftitle={Fractional-Order za SOFCs},
    pdfauthor={Your Name}
}
\usepackage[nameinlink,capitalise,noabbrev]{cleveref}

% ---------- Headers and footers ----------
\usepackage{fancyhdr}
\pagestyle{fancy}
\fancyhf{}
\fancyhead[L]{FODE equation discovery}
\fancyhead[R]{\nouppercase{\leftmark}}
\fancyfoot[C]{\thepage}
\setlength{\headheight}{14pt}

% ---------- Section formatting ----------
\usepackage{titlesec}
\titleformat{\chapter}[hang]
    {\normalfont\huge\bfseries}
    {\thechapter.}
    {0.6em}
    {}

% ---------- Document metadata ----------
\title{
    \textbf{Fractional-Order Differential Equation Modelling\\
    of Solid Oxide Fuel Cells}
}
\author{
    Mitja Mandić\\
    \small IJS E8
}
\date{\today}

\begin{document}

\section{Introduction}

In the following paper we are interested in fitting equivalent circuits to electrical impedance spectroscopy data in frequency domain with the help of formal grammars. We believe this
to be interesting, as it provides a straightforward way to exhaustively search the space of possible circuits up to a certain complexity. This offers flexibility not only
in the sense of fitting parameters, but also when it comes to building the underlying circuit.

\subsection{Solid Oxide Fuel Cells (SOFCs) and batteries}

Alternative sources of electrical energy are becoming increasingly relevant with the strong focus on shifting away from fossil fuel technology. In the present work we focus on data obtained
from two important options in the energy landscape - fuel cells, specifically SOFCs and batteries (Li-ion). In this section we briefly explain their working mechanics.

\subsubsection{Solid Oxide Fuel Cells}

Fuel cells can be viewed as factories for producing electricity. In a typical hydrogen fuel cell, hydrogen splits into protons and electron on the anode. The electrons
flow through an external circuit, while the protons cross the electrolyte layer and react with oxygen on the cathode, producing water and heat. SOFCs work differently; 
their high operating temperature enables the oxygen ions to move through a solid ceramic electrolyte from the cathode to the anode. There they react with the fuel such as hydrogen or carbon monoxide,
releasing electrons and heat and producing water or carbon dioxide, respectively.

Compared to standard fuel cells, SOFCs offer larger fuel flexibility and higher efficiency, while their high operating temperature also makes the waste heat useful, if handled correctly. Their main
drawbacks are slow startup and component wear. \cite{OHayre2016FuelCellFundamentals}

\subsubsection{Lithium Ion Batteries}

Lithium-ion batteries store energy chemically and release it as electricity. During discharge, lithium ions move from the negative electrode through the electrolyte to the positive electrode, 
while electrons flow through an external circuit to power a device. Charging drives both flows in the opposite direction, allowing the battery to be used again. 
Their main advantages are high energy density and rechargeability, though their capacity gradually declines with use.

\subsection{Electrochemical impedance spectroscopy (EIS) and equivalent circuit modeling (ECM)}

%explain what EIS data is briefly and why we use ECM
Electrochemical impedance spectroscopy is a modelling technique where we aim to understand the processes that occur in the battery or fuel cell better. In simple terms, we 
keep the cell at a certain potential, perturb the voltage slightly at a given frequency and measure the current's response. We then repeat this measurement at many frequencies, since the
electrochemical components behave differently at various frequencies.

Equivalent circuit modeling (ECM) looks at the measurement from a purely mathematical point of view, in the way that it is not concerned with the physical component from which it was obtained.
It aims to create a circuit from a predefined set of mostly simple elements that reproduces the measurements, in the hope of gaining some more understanding in the chemical process occurring
in the device that produced the data. 

\section{Grammar based modeling}

%Grammar based is good because it gives structure and it guarantees the circuits are ok. We can also give more rules
%based on domain knowledge

We use a formal context-free grammar to define candidate circuits which we then fit on the data. The grammar we use does not impose any structural
assumptions on the circuit. Each element can appear at any position, the only restriction being that a resistor $\term{Rs}$ is in series
with the rest of the circuit. The circuit elements present are standard (ECM) elements, which we present in Section \ref{circuit_gen} below.

\begin{equation}
\label{eq:grammar-relaxed}
\begin{aligned}
\nt{Circuit}
    &\longrightarrow
    \term{Rs}
    \mid
    \term{Rs}\;\term{+}\;\nt{Network},
    \\[2mm]
\nt{Network}
    &\longrightarrow
    \nt{Element}
    \\%\mid
    &\longrightarrow\;\nt{Network}\;\term{+}\;\nt{Network}\;
    %&\longrightarrow\term{(}\;\nt{Network}\;\term{+}\;\nt{Network}\;\term{)}
    \\
    &\longrightarrow%\quad%\mid
    \;\nt{Network}\;\term{||}\;\nt{Network}\;,
    %\term{(}\;\nt{Network}\;\term{||}\;\nt{Network}\;\term{)},
    \\[2mm]
\nt{Element}
    &\longrightarrow
    \term{R}
    \mid \term{L}
    \mid \term{CPE}
    \mid \term{W}.
\end{aligned}
\end{equation}

\section{Methodology}

Several steps were taken in circuit generation, subsequent derivation of impedance equations and finally parameter fitting.
Here we describe the core procedure.

\subsection{Circuit generation}\label{circuit_gen}

Circuit generation from formal grammars feels almost like solving a puzzle backwards. We consequently connect elements in series or parallel
and in the end map each $\nt{Element}$ to either a resistor, inductor, Warburg element or a CPE. We replace two resistors (inductors) in series or in parallel with a single resistor (inductor),
which helps to keep the total number of circuits down. This simplification is performed during generation to limit the computation time.

When generating circuits from grammars one has to define a certain complexity metric for each circuit. We decided to limit the number of elements in the circuit, as this more intuitively 
tells us how complex the circuit is, compared to, for example, grammar tree depth. We outline the generation algorithm for circuits up to a certain depth $N$ in Algorithm \ref{circuitGenAlgo}.
Function \texttt{Normalize} merges collapsible elements, namely resistors and inductors. Each circuit is also assumed to contain one resistor in series with the remainder.

\begin{algorithm}[H]\label{circuitGenAlgo}
\caption{Generate circuits with at most $N$ elements}
\begin{algorithmic}[1]
\Require Maximum element count $N$; allowed primitive elements or blocks $P$
\State $B_1,\ldots,B_{N-1} \gets \varnothing$; $C \gets \{\term{Rs}\}$ \Comment{Each circuit contains $\term{Rs}$}
\ForAll{$p \in P$}
    \State $q \gets \Call{Normalize}{p}$
    \If{$|q| < N$} \State $B_{|q|} \gets B_{|q|} \cup \{q\}$ \Comment{Normalized primitive blocks are building buckets of size 1} \EndIf 
\EndFor
\For{$n = 1,\ldots,N-1$}
    \For{$l = 1,\ldots,\lfloor n/2\rfloor$} \Comment{Circuit with $n$ elements is split in two sub-circuits with $l$ and $r$ elements}
        \State $r \gets n-l$
        \ForAll{$(u,v) \in B_l \times B_r$}
            \ForAll{$o \in \{\text{series},\text{parallel}\}$}
                \State $w \gets \Call{Normalize}{o(u,v)}$ \Comment{Normalize two connected already normalized buckets}
                \If{$|w| = n$} \State $B_n \gets B_n \cup \{w\}$ \Comment{If size of connected buckets is smaller than $n$ we skip it as it already exists}
                \EndIf 
            \EndFor
        \EndFor
    \EndFor
    \ForAll{$w \in B_n$}
        \State $C \gets C \cup \{\Call{Normalize}{\text{series}(\term{Rs},w)}\}$
    \EndFor
\EndFor
\State \Return $C$
\end{algorithmic}
\end{algorithm}

%Two resistors (inductors) in series or in parallel can be replaced by a single resistor (inductor),
%which
%
%However, we can take a few measures to 
%keep the number of circuit parameters as small as possible, namely:
%\begin{itemize}
%    \item replacing two resistors in series or in parallel with a single resistor
%    \item replacing two inductors in series or in parallel with a single inductor
%\end{itemize}

\subsection{Impedance equation generation}

Obtaining impedance equations from circuits generated by the grammar is straightforward. For series of elements we concatenate their
respective impedances, and for elements in parallel we sum their admittances. Below we list the elements present in the grammar and their impedances (admittance is the inverse of impedance).

\subsubsection*{Constant phase element (CPE)}
Constant phase element is a component that models an imperfect capacitor. It contains a dimensionless parameter $0\le\alpha\leq1$ that can be
interpreted as a measure of non-ideality. When $\alpha=1$ a CPE is a perfect capacitor.
CPE's impedance is
\[
    Z_{CPE} = \frac{1}{(j\omega)^{\alpha}Q}
\]
where $j$ is the imaginary unit, and $Q$ is a parameter independent of the frequency. $\omega=2\pi f$ is the angular frequency
\subsubsection*{Inductor (L) and Resistor (R)}

Resistor's impedance is constant does not depend on the frequency, and inductor's impedance is a linear function of the frequency. 
$$Z_R = R,$$
where $R$ is the resistor's resistance, and
\[
    Z_L = j\omega L,
\]
with $L$ the inductance.
\subsubsection*{Warburg element (W)}
Warburg element is used mainly in batteries EIS analysis. It is a special case of the CPE where $\alpha=0.5$. We focused on the infinite Warburg
element with impedance defined as
\[
    Z_W = \frac{\sigma}{\sqrt{jw}}
\]
\subsection{Parameter fitting}
As described above, frequency domain impedance equations are directly generated from the circuits by simply summing up elements
impedance formulas when they are connected in series, and summing up their reciprocal values when they are connected in parallel circuits.

Before we start the fitting procedure, we need to decide how complex circuits (i.e. how many elements they contain) are we fitting. Number of
candidates increases rapidly, which needs to be taken into account: there are 92 4 element circuits, 601 5 element circuits and several thousand 6 element
circuits. In the following work we focus on circuits up to 5 elements, as we deem they strike a balance between modest computation time and sufficient
accuracy.

Search for optimal parameter values can be split in two parts. In the first part we test all candidates, then select the best performing ones and 
continue optimizing the result.
%We run the optimization algorithm (more on that in a later section)
%up to a certain threshold, and for the circuits within a certain cutoff margin we re-run the optimization algorithm with 
%stricter thresholds to obtain results as close to full convergence as possible. 


\subsubsection{Initial value selection}

We opted for a multi-start stratified approach. Each circuit gets fitted 6 times in the initial evaluation. First fit is using deterministic
fixed values $p_{fixed}$:

\begin{itemize}
    \item 0.8 for $\alpha$
    \item 0.05 for $R_s$
    \item 0.02 for other resistors $R$
    \item 1.0 for $Q$
    \item 0.01 for $\sigma$
    \item $10^{-6}$ for $L$
\end{itemize}

For the remaining 5 starts we selected the values as follows:
\begin{itemize}
    \item $\alpha$: Divide interval $[0.03,0.98]$ in 5 equal parts and sample from each of the subintervals
    \item $R_s$: Divide interval $[0.05\cdot10^{-1},0.05\cdot10]=[0.005,0.5]$ into 5 parts and sample from subintervals
    \item For the other parameters: Divide $[p_{fixed}\cdot10^{-2},p_{fixed}\cdot10^{2}]$ in 5 subintervals and sample from them
\end{itemize}

This way we get a decent spread of initial values. We chose a narrower band for $R_s$ as it is more directly constrained by the higher frequencies,
and kept a broader band to allow for more uncertainty in initial values.

\subsubsection{Normalising residuals}\label{normalisation}

Residual normalisation is a standard procedure in robust statistics where we expect a certain amount of noise in the data.
To prevent large impedance measurements from dominating the result we decided to normalise the residuals during the optimization algorithm.
We tested three ways of normalising the residuals which we present below. 

Let \(Z_i\) be the measured impedance at frequency \(f_i\), and let
\(\widehat Z_i(\boldsymbol{\theta})\) be the circuit prediction. Define
the complex residual and its two real-valued components by
\begin{align}
r_i(\boldsymbol{\theta})
    &= \widehat Z_i(\boldsymbol{\theta})-Z_i, \\
r_{i,\mathrm{Re}}(\boldsymbol{\theta})
    &= \Re\!\left[r_i(\boldsymbol{\theta})\right], \\
r_{i,\mathrm{Im}}(\boldsymbol{\theta})
    &= \Im\!\left[r_i(\boldsymbol{\theta})\right].
\end{align}
The scaled residual supplied to the optimizer is
\begin{equation}
\widetilde r_{i,c}(\boldsymbol{\theta})
    = \frac{r_{i,c}(\boldsymbol{\theta})}{s_{i,c}},
    \qquad c\in\{\mathrm{Re},\mathrm{Im}\},
\end{equation}
where \(s_{i,c}>0\) is the scale for component \(c\) at frequency \(f_i\).

We tested the following choices of scale \(s_{i,c}\):
\begin{align}
s_{i,c}^{\mathrm{unscaled}}
    &= 1,\\
s_{i,c}^{\mathrm{relative}}
    &= \max\!\left(
        |[Z_i]_c|,
        0.1\,\operatorname{median}_{j}|[Z_j]_c|
       \right),\\
s_{i,\mathrm{Re}}^{\mathrm{magnitude}}
=s_{i,\mathrm{Im}}^{\mathrm{magnitude}}
    &= \max\!\left(
        |Z_i|,
        0.1\,\operatorname{median}_{j:\,|Z_j|>0}|Z_j|
       \right).
\end{align}

For spectrum data with repeated measurements \(Z_{ik}\) at frequency \(f_i\),
the scatter scale is calculated separately for each component:
\begin{align}
m_{i,c}
    &= \operatorname{median}_{k}[Z_{ik}]_c,\\
a_{i,c}
    &= 1.4826\,
       \operatorname{median}_{k}
       \left|[Z_{ik}]_c-m_{i,c}\right|,\\
s_{i,c}^{\mathrm{scatter}}
    &= \max\!\left(
        a_{i,c},
        \eta\,\operatorname{median}_{j}a_{j,c}
       \right),
\end{align}
where \(\eta\) is the scatter-floor fraction, which is in place to avoid division with very small numbers.



\subsubsection{Numerical algorithm}\label{numalg}

We fit parameters using several versions of least squares algorithms implemented in library SciPy \cite{2020SciPy-NMeth}. We test
the non-robust version which uses the ``linear'' loss function (standard least squares formulation),
and a version of the algorithm with the \texttt{soft\_l1} loss function in Equation \ref{soft_l1},

%In order to improve results of the algorithm we tried to improve robustness of SciPy's \texttt{least\_squares} algorithm. 
%For the initial optimization run we use:
%
%\begin{table}[]
%\centering
%\begin{tabular}{l|l}
%Parameter & value \\ \hline
%\texttt{x\_scale} & \texttt{jac} \\
%\texttt{ftol} &  10{-6}\\
%\texttt{xtol} &  10{-6}\\
%\texttt{gtol} &  10{-6}\\
%\texttt{f\_scale} &  0.001\\
%\texttt{max\_nfev} &  100\\
%\texttt{loss} &  \texttt{soft\_l1}\\
%\texttt{method} & \texttt{trf}
%\end{tabular}
%\label{parsopt1}
%\caption{Parameters for the SciPy.optimize least\_squares algorithm in the first optimization}
%\end{table}

%TRF method means we use the \textit{Trust region field} method of nonlinear least squares, and JAC scale is a specific 
%way of calculating matrix column scales. Tolerance fields (f-,x-,g-) stop the iteration if function values, parameters or function's derivatives
%move less than the threshold. Soft L1 loss with corresponding f-scale parameter means that instead of minimizing squared residuals, the optimizer
%is instead minimizing the following objective function:


\begin{equation}\label{soft_l1}
    F(\boldsymbol{x})
    = \frac{C^2}{2}
      \sum_{i=1}^{m}
      \rho\left(\frac{\widetilde r_i(\boldsymbol{x})^2}{C^2}\right),
\end{equation}
where $r_i(\boldsymbol{x})$ denotes the $i$th residual and
\begin{equation}
    \rho(z) = 2\left(\sqrt{1+z}-1\right).
\end{equation}

Parameter $C$ is called \texttt{f-scale} and is user defined. Residuals that are smaller than $C$
behave quadratically, as in standard formulation, but those that are larger only have an absolute effect
on the cost function - parameter $C$ intuitively presents a cutoff between quadratic and L1 loss.

The parameter estimations were obtained via two-stage optimization procedure.
We tested all candidate circuits in the first stage until they either stopped improving or
they were evaluated 100 times. For the second stage we kept those, that were within 20\%
of best performing one according optimizer cost function value at stopping point. We continue evaluating these circuits until they they reach a plateau and
perform final sorting via a measure, that takes model complexity into account.


\subsection{Model selection}

We attempted model selection in our experiments in two distinct ways. Firstly, we
robustified the well known BIC approach by remodelling the way quality of fit was
calculated. We call this approach \texttt{lp\_BIC}.

Our \texttt{combinatorial score} is an alternative approach to ranking models 
which takes into account circuit's structural complexity, not only the number
of parameters like in the \texttt{lp\_BIC} case.

\subsubsection{BIC based approaches}

%Residual scaling was used during parameter estimation but not in the final model-selection score. 
%These two stages serve different purposes: scaling improves the numerical robustness of the optimization 
%and balances contributions from different parts of the spectrum, whereas evaluating the fitted models in the 
%original impedance units provides a common criterion across normalization methods. 
%Using scaled residuals in the final score would make its value dependent on the chosen 
%normalization and would therefore complicate comparisons between scaling approaches. 
%To limit the influence of large raw residuals, we nevertheless replaced the 
%squared-error term with an Lₚ-type error using p=0.

Standard definition of BIC includes the MSE. Since we we normalise the residuals as
we fit the parameters, we decided to consistently alter the error 

use the \texttt{lp\_BIC}, where instead of raw 
we use weighted residuals to compute BIC.

Let \(N\) be the number of measured frequencies and \(k\) the number of
fitted circuit parameters. For each frequency \(i\), define

The post-fit error score is
\[
E_p=\frac{1}{N}\sum_{i=1}^{N}
\left(
|r_{i,\mathrm{Re}}|^p+|r_{i,\mathrm{Im}}|^p
\right),
\qquad p=0.5.
\]
We rank circuits using
\[
\mathrm{lp\_BIC}=2N\log(E_{0.5})+k\log(2N).
\]

\subsubsection{Combinatorial penalty}

The parameter-count penalty in BIC does not distinguish between circuits with
the same number of fitted parameters but different numbers of possible element
arrangements. The motivation behind this metric is that some circuit topologies
are simpler than others even if they contain the same number of elements.

To account for this additional structural flexibility, we define
the scaled post-fit error of a candidate circuit \(\mathcal{C}\) as

\begin{equation}\label{eq:scaled_lp_error}
\widetilde E_{0.5}(\mathcal{C})
=\frac{1}{N}\sum_{i=1}^{N}
\left(
\left|\widetilde r_{i,\mathrm{Re}}
    (\widehat{\boldsymbol{\theta}}_{\mathcal{C}})\right|^{0.5}
+
\left|\widetilde r_{i,\mathrm{Im}}
    (\widehat{\boldsymbol{\theta}}_{\mathcal{C}})\right|^{0.5}
\right),
\end{equation}
where \(\widehat{\boldsymbol{\theta}}_{\mathcal{C}}\) denotes the fitted
parameter vector for \(\mathcal{C}\). The combinatorial ranking score is
\begin{equation}\label{eq:combinatorial_score}
S_{\mathrm{comb}}(\mathcal{C})
=2N\log\!\left(\widetilde E_{0.5}(\mathcal{C})\right)
+\lambda_E e_{\mathcal{C}}
+\lambda_A\log M_{\mathcal{C}},
\end{equation}
where \(e_{\mathcal{C}}\) is the total number of circuit elements,
including \(R_s\), and \(M_{\mathcal{C}}\) is the number of distinct ways to
place these element types in the normalised series--parallel
topology of \(\mathcal{C}\). The series resistance remains fixed when
\(M_{\mathcal{C}}\) is counted. Permutations of identical elements and changes
in the order of series or parallel children do not produce new arrangements,
and assignments that simplify to a different topology are excluded.

The weights \(\lambda_E\) and \(\lambda_A\) control the penalties for adding
elements and for increasing the number of admissible arrangements,
respectively. In the current implementation, \(\lambda_E=9\) and
\(\lambda_A=0.5\). Circuits are ordered from the smallest to the largest value
of \(S_{\mathrm{comb}}\). 

%This is a custom ranking criterion and is not a BIC,
%because its structural penalties are chosen independently of a likelihood
%derivation.

%\subsection{}


%Simply put, more residuals behave quadratically with a larger f-scale, while a smaller f-scale means most of the function is behaving like
%the L1 loss. For the first optimization round, so for all circuits generated, the stopping point was the 100 function evaluations.
%
%Substituting these expressions gives the optimization problem
%\begin{equation}
%    \boxed{
%        \min_{\boldsymbol{x}} F(\boldsymbol{x})
%        =
%        \min_{\boldsymbol{x}}
%        10^{-6}\sum_{i=1}^{m}
%        \left[
%            \sqrt{1+10^{6}r_i(\boldsymbol{x})^2}-1
%        \right].
%    }
%\end{equation}
%
%In the second optimization round, that is, for 5 circuits with smallest MSE values, we again used the same algorithm, but with stricter
%optimization values presented in Table \ref{parsopt2}. Main difference is the increased number of steps and stricter function evaluation.
%In this case, for each of the circuits one of the tolerances was reached before we hit the step limit.
%\begin{table}[]
%\centering
%\begin{tabular}{l|l}
%Parameter & value \\ \hline
%\texttt{x\_scale} & \texttt{jac} \\
%\texttt{ftol} &  10{-8}\\
%\texttt{xtol} &  10{-8}\\
%\texttt{gtol} &  10{-8}\\
%\texttt{f\_scale} &  0.001\\
%\texttt{max\_nfev} &  10000\\
%\texttt{loss} &  \texttt{soft\_l1}\\
%\texttt{method} & \texttt{trf}
%\end{tabular}
%\caption{Parameters for the SciPy.optimize least\_squares algorithm in the second optimization round}
%\label{parsopt2}
%\end{table}


\section{Results}

We apply methods described above to data from both solid oxide fuel cells, as well as batteries. We try to distinguish which of the numerical adaptations
of the algorithm (additional settings for the SciPy's least-squares solver, normalising residuals) have the largest impact on the fit.

The plots presented in the graphs are the Nyquist graphs of the 5 circuits which we fitted until convergence, ordered by their BIC values.

\subsection{Simulation study}

We first test the above method on two simulated datasets, one coming from an ECM that is typical for SOFCs, and the second reproducing an impedance
dataset that resembles a battery. 

\[
\begin{aligned}
\textnormal{SOFC circuit}
    &= R_s
    \seriesconn
    \left[
        R_1
        \parallel
        \CPEel{Q_2}{\alpha_2}
    \right]
    \seriesconn
    \left[
        R_3
        \parallel
        \CPEel{Q_4}{\alpha_4}
    \right]
    ,
    \\[1.5ex]
\textnormal{Battery circuit}
    &= R_s
    \seriesconn
    \left[
        R_1
        \parallel
        \CPEel{Q_2}{\alpha_2}
    \right]
    \seriesconn
    W(\sigma_3).
\end{aligned}
\]

We simulated the spectra by evaluating the above circuits at 351 logarithmically spaced points between 0.1 Hz and 100 kHz. At each frequency point we
simulated 20 measurements by adding independent Gaussian with standard deviation $\sqrt{(0.15|Z|^2) + (0.05\Omega)^2}$ to both
real and imaginary component.

\begin{figure}[htbp]
    \begin{subfigure}[t]{0.48\textwidth}
        \centering
        \includegraphics[width=\linewidth]{/home/mitek/Work/Equation-Discovery-for-FODE/Results/campaigns/simulated_006_007_no_G_relaxed_5_relative_soft_l1_fscale_0p001_optimizer_value_lp_bic_1start/runs/simulated_simulated_spectrum_006/Results/Plots/Nyquists/Fits of simulated battery data.png}
        \caption{Simulated battery data fit results}
        \label{batterySim}
    \end{subfigure}
    \hfill
    \begin{subfigure}[t]{0.48\textwidth}
        \centering
        \includegraphics[width=\linewidth]{/home/mitek/Work/Equation-Discovery-for-FODE/Results/campaigns/simulated_006_007_no_G_relaxed_5_relative_soft_l1_fscale_0p001_optimizer_value_lp_bic_1start/runs/simulated_simulated_spectrum_007/Results/Plots/Nyquists/Fits of simulated SOFC data.png}
        \caption{Simulated SOFC data fit results}
        \label{sofcSim}
    \end{subfigure}
    \caption{Comparison of the two fits.}
    \label{simData}
\end{figure}

In the Figure \ref{simData} we see that even with substantial noise, our grammar based method manages to recreate the Nyquist plots. Looking especially
at Figure \ref{sofcSim} it becomes clear that several circuits are capable of recreating the data. The selected circuit exactly matches the simulating one 
in the battery data case, which is not the case in the SOFC simulated data, where a completely different circuit structure matches the Nyquist plot. This already
raises the question of identifiability.


\[
\textnormal{Selected SOFC circuit}
    = R_s
    \seriesconn
    \left[
    \CPEel{Q_1}{\alpha_1}
    \parallel
    \left(
        R_2
        \seriesconn
        \left(
            R_3
            \parallel
            \CPEel{Q_4}{\alpha_4}
        \right)
    \right)
    \right]
\]

\begin{table}[htbp]
\centering
\caption{True and selected parameters for simulated battery spectrum}
\label{tab:sim006-parameters}
\begin{tabular}{@{}lrlr@{}}
\toprule
\textbf{True parameter} & \textbf{Value}
& \textbf{Selected parameter} & \textbf{Estimate} \\
\midrule
$R_s$       & 0.100000 & $R_s$         & 0.098178 \\
$R_1$       & 0.400000 & $R_1$         & 0.400238 \\
$Q_2$       & 0.020000 & $Q_2$         & 0.020729 \\
$\alpha_2$  & 0.850000 & $\alpha_2$    & 0.852767 \\
$\sigma_3$    & 0.500000 & $\sigma_3$    & 0.497347 \\
\bottomrule
\end{tabular}
\end{table}

In Table \ref{tab:sim006-parameters} we see that our method successfully found the correct model with parameter estimation that very closely
matches the simulator values.

%\begin{table}[htbp]
%\centering
%\caption{True and selected parameters for simulated SOFC spectrum 007.
%The selected circuit has a different topology, so values placed on the
%same row should not be interpreted as estimates of the same component.}
%\label{tab:sim007-parameters}
%\begin{tabular}{@{}lrlr@{}}
%\toprule
%\textbf{True parameter} & \textbf{Value}
%& \textbf{Selected parameter} & \textbf{Estimate} \\
%\midrule
%$R_s$       & 0.100000 & $R_s$       & 0.099206 \\
%$R_1$       & 0.400000 & $R_2$       & 0.407632 \\
%$Q_1$       & 0.020000 & $Q_1$       & 0.019695 \\
%$\alpha_1$  & 0.850000 & $\alpha_1$  & 0.860670 \\
%$R_2$       & 0.800000 & $R_3$       & 0.698284 \\
%$Q_2$       & 5.000000 & $Q_4$       & 5.625948 \\
%$\alpha_2$  & 0.700000 & $\alpha_4$  & 0.771703 \\
%\bottomrule
%\end{tabular}
%\end{table}

\subsection{SOFC data}
We continue testing our algorithm on real world data, first from a solid oxide fuel cell and in the following Section \ref{battery} we present results from battery data.
For SOFC we chose \texttt{eis\_20200226\_230006.npz} dataset from repository \\\href{https://portal.ijs.si/nextcloud/s/xTa2cmtfxXn2jSz}{https://portal.ijs.si/nextcloud/s/xTa2cmtfxXn2jSz}.
In this dataset, 233 measurements are repeated several times at each frequency and we calculate their median to obtain one impedance value per frequency.

In Figure \ref{sofc4el} we see results of fitting circuits with up to 4 elements on the SOFC dataset, which delivers poor results as the details of the data are not 
captured in the fit.

\begin{figure}[h]
\centering
\includegraphics[width=\textwidth]{../../Results/campaigns/unified_no_G_relaxed_4_elements/runs/spectrum_eis_20200226_230006/Results/Plots/Nyquists/eis_20200226_230006_no_G_relaxed_4_best_5_refined.png}
\caption{Nyquist graphs of 5 best fitting circuits on SOFC dataset with up to 4 elements after MSE margin reoptimization}
\label{sofc4el}
\end{figure}

Including more elements in the search by itself does not improve the result visually as reported in Figure
\ref{sofc5el_nonorm}. The fit is closer to the data especially in higher frequency range on the left hand side,
but still fails to capture the general shape of data.

Including relative normalisation brings improvement to the fit, which
we present in Figure \ref{sofc5el_relScale}. We now have a clear dip around 0.057 on the
real axis, and low frequency shape on the right of the spectrum is captured closely. 

We can further improve the fit by leveraging data abundance in this case, where measurements of each 
fuel cell were repeated across frequencies. We present these results in Figure \ref{sofc5el_scatScale},
where we further improve the high frequency fit. It has to be noted, that even though the circuit appears
to be fitting closer to the data, the same structure is selected in both relative scaling case and
scatter scaling case. 

Parameter values for both fits are presented in Table \ref{tab:sofc_relative_scatter_parameters}. 
First three listed parameters are nearly identical, and in both circuits $\alpha_2$ falls on the upper 
boundary of the feasible interval, making the CPE in this case behave nearly like a perfect capacitor. 
There are larger differences in parameter values for the final resistor and CPE. 

Looking at the 
residual plot (INSERT RESIDUAL PLOTS) we can see the largest variation of estimation quality at the both ends of the
observed spectrum. Scatter scaling uses information from repeated measurements, therefore the scales obtained from
more datasources can help reduce the influences of high measurement variability.

For \texttt{f\_scale} parameter $0.001$, almost all residuals fall beneath that mark. This makes the loss function behave effectively like
an L1 loss function, but we choose to keep the \texttt{soft\_l1} formulation to preserve the differentiability and numerical stability provided
by the SciPy implementation.

\begin{figure}[h]
\centering
\includegraphics[width=\textwidth]{../../Results/campaigns/unified_no_G_relaxed_5_elements_unscaled/runs/spectrum_eis_20200226_230006/Results/Plots/Nyquists/refined_scaled_lp_bic.png}
\caption{Nyquist graphs of 5 best fitting circuits on SOFC dataset with up to 4 elements after MSE margin reoptimization}
\label{sofc5el_nonorm}
\end{figure}

\begin{figure}[h]
\centering
\includegraphics[width=\textwidth]{../../Results/campaigns/real_pcl_spectrum_no_G_relaxed_5_relative_soft_l1_fscale_0p001_optimizer_value_lp_bic_6starts_20261112/runs/spectrum_eis_20200226_230006/Results/Plots/Nyquists/eis_20200226_230006_no_G_relaxed_5_best_5_refined.png}
\caption{SOFC 5 el rel scale}
\label{sofc5el_relScale}
\end{figure}

\begin{figure}[h]
\centering
\includegraphics[width=\textwidth]{../../Results/campaigns/real_pcl_sofc_no_G_relaxed_5_soft_l1_scatter_fscale_0p001_optimizer_value_lp_bic_6starts_20261113/runs/spectrum_eis_20200226_230006/Results/Plots/Nyquists/eis_20200226_230006_no_G_relaxed_5_best_5_refined.png}
\caption{SOFC 5 el scat scale}
\label{sofc5el_scatScale}
\end{figure}



\begin{table}[htbp]
\centering
\caption{Comparison of the fitted SOFC circuit parameters obtained with
relative and scatter residual scaling.}
\label{tab:sofc_relative_scatter_parameters}
\begin{tabular}{lrr}
\toprule
Parameter & Relative scaling & Scatter scaling \\
\midrule
\(R_s\) (\si{\ohm})
    & 0.0474021 & 0.0463463 \\

\(R_1\) (\si{\ohm})
    & 0.0241604 & 0.0255941 \\

\(Q_2\) (\si{\siemens\second\tothe{\alpha_2}})
    & 1.15428 & 1.09144 \\

\(\alpha_2\)
    & 0.990000 & 0.990000 \\

\(R_3\) (\si{\ohm})
    & 0.0127951 & 0.0176448 \\

\(Q_4\) (\si{\siemens\second\tothe{\alpha_4}})
    & 0.308221 & 0.682791 \\

\(\alpha_4\)
    & 0.958089 & 0.778419 \\
\bottomrule
\end{tabular}
\end{table}

\begin{table}[htbp]
\centering
\caption{SOFC fit comparison for different residual scaling methods.}
\label{tab:sofc_scaling_comparison}
\begin{tabular}{lrrr}
\toprule
Metric & No normalisation & Relative scaling & Scatter scaling \\
\midrule
MSE ($\Omega^2$)       & $2.6974\times10^{-6}$ & $3.2739\times10^{-6}$ & $3.1278\times10^{-6}$ \\
BIC                    & $-5932.61$ & $-5842.35$ & $-5863.62$ \\
$\mathrm{lp\_BIC}$     & $-1402.54$ & $-1374.24$ & $-1453.17$ \\
\bottomrule
\end{tabular}
\end{table}



%\begin{figure}[h!]
%\centering
%\includegraphics[width=\textwidth]{../../Results/Plots/Nyquists/bulk/eis_20200227_070006_no_G_relaxed_5_best_5_fixed_starting_values_no_normalization.png}
%\caption{Nyquist graphs of 5 best fitting circuits according to our algorithm with up to 5 elements}
%\label{sofc5el_nonorm}
%\end{figure}
%
%Before expanding the search space with even more circuits, we decided to try to make the algorithm more robust.
%We first implemented normalising residuals at each step of the algorithm and present results in Figure \ref{sofc5el}. There is noticable
%improvement in how closely the graphs fit the data. When inspecting the results in Table \ref{fig:first-table-circuits} more closely, we see that
%several parameters are close to their limits (especially $\alpha$). 
%
%\begin{figure}[h!]
%\centering
%\includegraphics[width=\textwidth]{../../Results/Plots/Nyquists/bulk/eis_20200227_070006_no_G_relaxed_5_best_5.png}
%\caption{Nyquist graphs of 5 best fitting circuits according to our algorithm with up to 5 elements and normalised residuals.}
%\label{sofc5el}
%\end{figure}
%
%To improve on unstable solutions we robustify the algorithm with adaptations described in Section \ref{numalg}. This way we get some more diversity
%in results in Graph \ref{sofc5elnum}.


\begin{table}[htbp]
\centering
\caption{SOFC fit comparison for different values of \(f_{\mathrm{scale}}\).}
\label{tab:sofc_fscale_comparison}
\begin{tabular}{lrrr}
\toprule
Metric
& \(f_{\mathrm{scale}}=0.001\)
& \(f_{\mathrm{scale}}=0.01\)
& \(f_{\mathrm{scale}}=0.1\) \\
\midrule
MSE (\(\Omega^2\))
& \(3.2739\times10^{-6}\)
& \(2.5932\times10^{-6}\)
& \(2.6254\times10^{-6}\) \\

BIC
& \(-5842.35\)
& \(-5950.97\)
& \(-5945.22\) \\

\(\mathrm{lp\_BIC}\)
& \(-1374.24\)
& \(-1339.74\)
& \(-1319.68\) \\

\(\mathrm{scaled\_lp\_BIC}\)
& \(-180.57\)
& \(-188.34\)
& \(-174.30\) \\
\bottomrule
\end{tabular}
\end{table}

\begin{figure}[htbp]
    \begin{subfigure}[t]{0.48\textwidth}
        \centering
        \includegraphics[width=\linewidth]{../../Results/campaigns/unified_no_G_relaxed_5_elements_relative_soft_l1_fscale_0p01/runs/spectrum_eis_20200226_230006/Results/Plots/Nyquists/eis_20200226_230006_no_G_relaxed_5_best_5_refined.png}
        \caption{Results with f-scale set to 0.01}
        \label{fscale_0.01comp}
    \end{subfigure}
    \hfill
    \begin{subfigure}[t]{0.48\textwidth}
        \centering
        \includegraphics[width=\linewidth]{../../Results/campaigns/unified_no_G_relaxed_5_elements_relative_soft_l1_fscale_0p001/runs/spectrum_eis_20200226_230006/Results/Plots/Nyquists/eis_20200226_230006_no_G_relaxed_5_best_5_refined.png}
        \caption{Results with f-scale set to 0.001}
        \label{fscale_0.001comp}
    \end{subfigure}
    \caption{Comparison between different f-scale values for the optimizer}
    \label{fscalecomp}
\end{figure}


\clearpage

\subsection{Battery data}\label{battery}

The \texttt{eis\_bank1\_20251016-084621.pcl} battery spectrum contains 128 complex impedance values for channel \texttt{batt\_1}, spanning 0.00141 to 1000 Hz.
The real part ranges from 0.0165 to 0.0392 $\Omega$. The imaginary part remains negative across the entire spectrum, reaching -0.0295 $\Omega$ at the lowest frequency. 
Both components increase in magnitude toward low frequencies, producing a pronounced low-frequency tail that is normal in batteries and implies low frequency
diffusion. In this case measurements are not repeated (different channels mean different physical battery packs) and we therefore do not take the
scatter scaling into account as we did for SOFC data.

In Figure \ref{batt_relScale} the 5 element fit is presented. After the first round, only one circuit was in the 20\% reoptimization range. In this case it is
more difficult to inspect the fit quality visually, as the measurements are more spread out on the imaginary axis. The prominent tail low frequency tail on the
right side corresponds to the slower ion diffusion taking place. 

\begin{figure}[h]
\centering
\includegraphics[width=\textwidth]{../../Results/campaigns/real_pcl_spectrum_no_G_relaxed_5_relative_soft_l1_fscale_0p001_optimizer_value_lp_bic_6starts_20261112/runs/pcl_eis_bank1_20251016-084621_batt_1/Results/Plots/Nyquists/eis_bank1_20251016-084621_batt_1_no_G_relaxed_5_best_5_refined.png}
\caption{Battery 5 el rel scale}
\label{batt_relScale}
\end{figure}


\begin{table}[htbp]
\centering
\begin{tabular}{ll}
\textbf{Metric} & \textbf{Value} \\ \hline
Element count & 5 \\
Parameter count & 7 \\
MSE & $1.49010\times10^{-7}$ \\
BIC & $-3985.31221$ \\ \hline
\textbf{Circuit parameter} \\\hline
$R_s$ (\si{\ohm}) & 0.0155126 \\
$R_3$ (\si{\ohm}) & 0.00593952 \\
$L_4$ (\si{\henry}) & 0.734682 \\
$Q_1$ (\si{\siemens\second\tothe{\alpha_1}}) & 2.91787 \\
$\alpha_1$ & 0.673854 \\
$Q_2$ (\si{\siemens\second\tothe{\alpha_2}}) & 433.414 \\
$\alpha_2$ & 0.602733
\end{tabular}
\caption{Fitted parameters of the best-performing circuit.}
\label{tab:best_circuit_sofc}
\end{table}



\begin{table}[htbp]
\centering
\caption{Battery fit comparison for different values of
\(f_{\mathrm{scale}}\).}
\label{tab:battery_fscale_comparison}
\begin{tabular}{lrrr}
\toprule
Metric
& \(f_{\mathrm{scale}}=0.001\)
& \(f_{\mathrm{scale}}=0.01\)
& \(f_{\mathrm{scale}}=0.1\) \\
\midrule
MSE (\(\Omega^2\))
& \(2.7692\times10^{-7}\)
& \(2.6202\times10^{-7}\)
& \(2.5071\times10^{-7}\) \\

BIC
& \(-3826.67\)
& \(-3840.83\)
& \(-3852.12\) \\

\(\mathrm{lp\_BIC}\)
& \(-966.10\)
& \(-953.29\)
& \(-940.70\) \\

\(\mathrm{scaled\_lp\_BIC}\)
& \(\mathbf{-300.82}\)
& \(-293.54\)
& \(-284.62\) \\
\bottomrule
\end{tabular}
\end{table}





\clearpage
\nocite{*} 
\printbibliography
%\bibliographystyle{plainnat}

%\bibliography{references} 

% TODO: Document the fitting, refinement, and ranking settings used for each
% reported run; distinguish the final method from exploratory combinations.

\end{document}
